\documentclass[12pt]{article}
\usepackage[spanish]{babel}
\usepackage[utf8]{inputenc}
\usepackage[T1]{fontenc}
\usepackage{amsmath, amssymb}
\usepackage{geometry}
\usepackage{enumitem}
\usepackage{needspace}
\usepackage{tikz}
\usetikzlibrary{babel}

\geometry{letterpaper, margin=2cm}
\setlength{\parindent}{0pt}
\setlength{\parskip}{6pt}

\newcommand{\opcion}[2]{\item[#1)] $#2$}
\newcommand{\puntografico}{0.35}

% Mini grafico con P en origen
% grafcirc{cx}{cy}{r}{etiq-centro}{etiq-eje-x}
% scale=0.078 → necesita puntos grandes (1.0) para que sean visibles
\newcommand{\grafcirc}[5]{%
\begin{tikzpicture}[scale=0.078]
  \path[use as bounding box] (-32,-15) rectangle (43,34);
  \draw[<->, thick] (-30,0) -- (36,0);
  \draw[<->, thick] (0,-12) -- (0,31);
  \node[anchor=west] at (37,-1.2) {\scriptsize$x$ (este)};
  \node[anchor=south east] at (7,32.5) {\scriptsize$y$ (norte)};
  \fill (0,0) circle (1.0) node[below left] {\scriptsize$P$};
  \draw[dashed] (#1,0) -- (#1,{#2+#3});
  \draw[dashed] (0,#2) -- (#1,#2);
  \draw[dashed] (0,{#2+#3}) -- (#1,{#2+#3});
  \draw[thick] (#1,#2) circle (#3);
  \fill (#1,#2) circle (1.0);
  \fill (#1,{#2+#3}) circle (1.0);
  \node[anchor=south west, inner sep=1pt, xshift=2pt, yshift=1pt] at (#1,#2) {\scriptsize$#4$};
  \node[anchor=north] at (#1,0) {\scriptsize$#5$};
  \node[left] at (0,#2) {\scriptsize$#2$};
  \node[left] at (0,{#2+#3}) {\scriptsize\pgfmathparse{int(#2+#3)}\pgfmathresult};
\end{tikzpicture}}

% Mini grafico sin P: nueva posicion tras traslacion
% grafpos{cx}{cy}{r}{etiq-centro}{etiq-eje-x}
\newcommand{\grafpos}[5]{%
\begin{tikzpicture}[scale=0.35]
  \path[use as bounding box] (-2,-1) rectangle (11,11);
  \draw[<->, thick] (-1,0) -- (10,0);
  \draw[<->, thick] (0,-0.5) -- (0,10);
  \node[anchor=west] at (10.2,-0.3) {\scriptsize$x$};
  \node[anchor=south east] at (1.1,10.2) {\scriptsize$y$};
  \draw[dashed] (#1,0) -- (#1,{#2+#3});
  \draw[dashed] (0,#2) -- (#1,#2);
  \draw[dashed] (0,{#2+#3}) -- (#1,{#2+#3});
  \draw[thick] (#1,#2) circle (#3);
  \fill (#1,#2) circle (0.22) node[anchor=south west, inner sep=1pt] {\scriptsize$#4$};
  \fill (#1,{#2+#3}) circle (0.22);
  \node[anchor=north] at (#1,0) {\scriptsize$#5$};
  \node[left] at (0,#2) {\scriptsize$#2$};
  \node[left] at (0,{#2+#3}) {\scriptsize\pgfmathparse{int(#2+#3)}\pgfmathresult};
\end{tikzpicture}}

% Variante con eje Y extendido a 12 (para cy+r=10)
\newcommand{\grafposX}[5]{%
\begin{tikzpicture}[scale=0.35]
  \path[use as bounding box] (-2,-1) rectangle (11,13);
  \draw[<->, thick] (-1,0) -- (10,0);
  \draw[<->, thick] (0,-0.5) -- (0,12);
  \node[anchor=west] at (10.2,-0.3) {\scriptsize$x$};
  \node[anchor=south east] at (1.1,12.2) {\scriptsize$y$};
  \draw[dashed] (#1,0) -- (#1,{#2+#3});
  \draw[dashed] (0,#2) -- (#1,#2);
  \draw[dashed] (0,{#2+#3}) -- (#1,{#2+#3});
  \draw[thick] (#1,#2) circle (#3);
  \fill (#1,#2) circle (0.22) node[anchor=south west, inner sep=1pt] {\scriptsize$#4$};
  \fill (#1,{#2+#3}) circle (0.22);
  \node[anchor=north] at (#1,0) {\scriptsize$#5$};
  \node[left] at (0,#2) {\scriptsize$#2$};
  \node[left] at (0,{#2+#3}) {\scriptsize\pgfmathparse{int(#2+#3)}\pgfmathresult};
\end{tikzpicture}}

\begin{document}

\begin{center}
  {\Large \textbf{Pr\'actica para la Prueba Nacional Estandarizada Sumativa 2026}}\\[4pt]
  {\large Secundaria -- Bloque~1 -- Afirmaciones~1~y~2}\\[4pt]
  Nombre: \rule{8cm}{0.4pt}\hfill Fecha: \rule{3cm}{0.4pt}
\end{center}

\medskip
\textbf{Instrucciones:} Seleccione la \'unica opci\'on correcta en cada \'item.
Tiempo disponible: \textbf{21 minutos}.

\begin{enumerate}[label=\textbf{\arabic*.}, itemsep=1.0cm]

%---------------------------------------
% ITEM 1
%---------------------------------------
\Needspace{0.82\textheight}
\item
Un dr\'on de monitoreo emite una se\~nal inal\'ambrica que llega a $8$ m a su alrededor.
La ubicaci\'on de ese dr\'on ($D$) es $20$ m al este y $12$ m al norte de un punto de
control ($P$), el cual se considera como origen.

\textquestiondown Cu\'al de las siguientes representaciones gr\'aficas, en la que las unidades
est\'an en metros, corresponde a la circunferencia del alcance m\'aximo de la se\~nal que
emite ese dr\'on?

\vspace{1.8\baselineskip}
\begin{center}
\setlength{\tabcolsep}{0pt}
\renewcommand{\arraystretch}{1}
\begin{tabular}{@{}r@{\hspace{0.45cm}}c@{\hspace{1.6cm}}r@{\hspace{0.45cm}}c@{}}
\raisebox{1.65cm}{\textbf{A)}} & \grafcirc{20}{12}{8}{D}{20} &
\raisebox{1.65cm}{\textbf{B)}} & \grafcirc{12}{20}{8}{D}{12} \\[0.85cm]
\raisebox{1.65cm}{\textbf{C)}} & \grafcirc{20}{12}{4}{D}{20} &
\raisebox{1.8cm}{\textbf{D)}} & \begin{tikzpicture}[scale=0.14]
  \path[use as bounding box] (-30,-9) rectangle (12,24);
  \draw[<->, thick] (-28,0) -- (5,0);
  \draw[<->, thick] (0,-7) -- (0,23);
  \node[anchor=west] at (5.3,-0.7) {\scriptsize$x$ (este)};
  \node[anchor=south east] at (3,23.5) {\scriptsize$y$ (norte)};
  \fill (0,0) circle (0.55) node[below left] {\scriptsize$P$};
  \draw[dashed] (-20,0) -- (-20,20);
  \draw[dashed] (0,12) -- (-20,12);
  \draw[dashed] (0,20) -- (-20,20);
  \draw[thick] (-20,12) circle (8);
  \fill (-20,12) circle (0.55);
  \fill (-20,20) circle (0.55);
  \node[anchor=south west, inner sep=1pt, xshift=2pt, yshift=1pt] at (-20,12) {\scriptsize$D$};
  \node[anchor=north] at (-20,0) {\scriptsize$-20$};
  \node[left] at (0,12) {\scriptsize$12$};
  \node[left] at (0,20) {\scriptsize$20$};
\end{tikzpicture}
\end{tabular}
\end{center}

%---------------------------------------
% ITEM 2
%---------------------------------------
\Needspace{0.72\textheight}
\item
La siguiente representaci\'on gr\'afica, en la que las unidades est\'an en metros, muestra
la ubicaci\'on del centro ($S$) de un sistema de riego circular y la circunferencia
correspondiente a su alcance.

\begin{center}
\begin{tikzpicture}[scale=0.34]
  \path[use as bounding box] (-3,-3) rectangle (20,16);
  \draw[<->, thick] (-2,0) -- (18,0);
  \draw[<->, thick] (0,-2) -- (0,15);
  \node[anchor=west] at (18.4,-0.8) {$x$ (este)};
  \node[anchor=south east] at (3.5,15.5) {$y$ (norte)};
  \draw[dashed] (10,0) -- (10,6);
  \draw[dashed] (0,6) -- (10,6);
  \draw[thick] (10,6) circle (6);
  \fill (10,6) circle (\puntografico) node[above left, inner sep=2pt] {$S$};
  \node[anchor=north] at (10,0) {$10$};
  \node[left] at (0,6) {$6$};
\end{tikzpicture}
\end{center}

De acuerdo con la informaci\'on anterior, \textquestiondown cu\'al es la representaci\'on
algebraica, en la que las unidades est\'an en metros, de la circunferencia correspondiente
al alcance de ese sistema de riego?

\begin{enumerate}[label=\Alph*)]
  \opcion{A}{(x+10)^{2}+(y+6)^{2}=36}
  \opcion{B}{(x-10)^{2}+(y-6)^{2}=36}
  \opcion{C}{(x-6)^{2}+(y-10)^{2}=36}
  \opcion{D}{(x-10)^{2}+(y-6)^{2}=6}
\end{enumerate}

%---------------------------------------
% ITEM 3
%---------------------------------------
\Needspace{0.82\textheight}
\item
Un juego consiste en lanzar una bolsa de arena hacia una zona circular en el suelo,
cuya medida del radio es $5$ m.

La siguiente representaci\'on gr\'afica, cuyas unidades est\'an en metros, muestra la
circunferencia correspondiente al borde de esa zona:

\begin{center}
\begin{tikzpicture}[scale=0.34]
  \path[use as bounding box] (-2,-2) rectangle (17,14);
  \draw[<->, thick] (-1,0) -- (16,0);
  \draw[<->, thick] (0,-1) -- (0,13);
  \node[anchor=west] at (16.2,-0.4) {$x$};
  \node[anchor=south east] at (1.2,13.3) {$y$};
  \draw[thick] (8,6) circle (5);
  \draw[dashed] (8,0) -- (8,11);
  \draw[dashed] (0,6) -- (8,6);
  \draw[dashed] (0,11) -- (8,11);
  \fill (8,6) circle (\puntografico);
  \fill (8,11) circle (\puntografico);
  \node[anchor=north] at (8,0) {$8$};
  \node[left] at (0,6) {$6$};
  \node[left] at (0,11) {$11$};
\end{tikzpicture}
\end{center}

Adem\'as, durante el juego:
\begin{itemize}
  \item Ana lanz\'o una bolsa y esta se ubic\'o en el punto correspondiente a $(3,\,7)$.
  \item Bruno lanz\'o una bolsa y esta se ubic\'o en el punto correspondiente a $(8,\,12)$.
  \item Carla lanz\'o una bolsa y esta se ubic\'o en el punto correspondiente a $(11,\,9)$.
  \item Diego lanz\'o una bolsa y esta se ubic\'o en el punto correspondiente a $(13,\,7)$.
\end{itemize}

De acuerdo con la informaci\'on anterior, \textquestiondown cu\'al persona lanz\'o la
bolsa que se ubic\'o en el interior de esa zona?

\begin{enumerate}[label=\Alph*)]
  \item Ana
  \item Bruno
  \item Carla
  \item Diego
\end{enumerate}

%---------------------------------------
% ITEM 4
%---------------------------------------
\Needspace{0.84\textheight}
\item
Una antena de telecomunicaciones ($T$) emite una se\~nal que permite a los tel\'efonos
celulares ubicados a $5$ km o menos alrededor de ella realizar y recibir llamadas.
La siguiente representaci\'on gr\'afica, en la que las unidades est\'an en kil\'ometros,
muestra la ubicaci\'on de la antena y la circunferencia correspondiente al alcance m\'aximo
de la se\~nal.

\begin{center}
\begin{tikzpicture}[scale=0.34]
  \path[use as bounding box] (-2,-2) rectangle (15,14);
  \draw[<->, thick] (-1,0) -- (14,0);
  \draw[<->, thick] (0,-1) -- (0,13);
  \node[anchor=west] at (14.4,-0.4) {$x$ (este)};
  \node[anchor=south east] at (1.3,13.4) {$y$ (norte)};
  \draw[thick] (7,5) circle (5);
  \draw[dashed] (7,0) -- (7,10);
  \draw[dashed] (0,5) -- (7,5);
  \draw[dashed] (0,10) -- (7,10);
  \fill (7,5) circle (\puntografico)
    node[anchor=south east, inner sep=2pt, xshift=-2pt, yshift=2pt] {$T$};
  \fill (7,10) circle (\puntografico);
  \node[anchor=north] at (7,0) {$7$};
  \node[left] at (0,5) {$5$};
  \node[left] at (0,10) {$10$};
\end{tikzpicture}
\end{center}

De acuerdo con la informaci\'on anterior, si las ubicaciones que tienen los tel\'efonos
celulares $P$, $Q$, $S$ y $U$ corresponden, respectivamente, a los puntos
$(10,\,9)$, $(13,\,5)$, $(7,\,11)$ y $(2,\,1)$, entonces \textquestiondown en cu\'al de esos
tel\'efonos se pueden realizar y recibir llamadas por medio de la se\~nal de esa antena?

\begin{enumerate}[label=\Alph*)]
  \item $P$
  \item $Q$
  \item $S$
  \item $U$
\end{enumerate}

%---------------------------------------
% ITEM 5
%---------------------------------------
\Needspace{0.84\textheight}
\item
Considere la siguiente representaci\'on gr\'afica, en la que las unidades est\'an en
kil\'ometros, de la circunferencia correspondiente al alcance m\'aximo de la se\~nal
emitida por una antena ubicada en el pueblo $H$\text{:}

\begin{center}
\begin{tikzpicture}[scale=0.45]
  \path[use as bounding box] (-1.5,-2) rectangle (9,9);
  \draw[<->, thick] (-1,0) -- (8.5,0);
  \draw[<->, thick] (0,-1) -- (0,8.5);
  \node[anchor=west] at (8.7,-0.3) {$x$ (este)};
  \node[anchor=south east] at (1.2,8.7) {$y$ (norte)};
  \draw[thick] (4,5) circle (2);
  \draw[dashed] (4,0) -- (4,7);
  \draw[dashed] (0,5) -- (4,5);
  \draw[dashed] (0,7) -- (4,7);
  \fill (4,5) circle (0.22)
    node[anchor=south east, inner sep=2pt] {$H$};
  \fill (4,7) circle (0.22);
  \node[anchor=north] at (4,0) {$4$};
  \node[left] at (0,5) {$5$};
  \node[left] at (0,7) {$7$};
\end{tikzpicture}
\end{center}

De acuerdo con la informaci\'on anterior, para mejorar el servicio la empresa traslad\'o la
antena al pueblo $K$, el cual se ubica a $3$ km al este y $3$ km al sur de $H$. \textquestiondown Cu\'al es
la representaci\'on gr\'afica, en la que las unidades est\'an en kil\'ometros, de la
circunferencia correspondiente al alcance m\'aximo de la se\~nal emitida por la antena en
la nueva ubicaci\'on?

\vspace{1.5\baselineskip}
\begin{center}
\setlength{\tabcolsep}{0pt}
\renewcommand{\arraystretch}{1}
\begin{tabular}{@{}r@{\hspace{0.4cm}}c@{\hspace{1.4cm}}r@{\hspace{0.4cm}}c@{}}
\raisebox{1.9cm}{\textbf{A)}} & \grafpos{7}{5}{2}{K}{7} &
\raisebox{1.9cm}{\textbf{B)}} & \grafposX{4}{8}{2}{K}{4} \\[0.6cm]
\raisebox{1.9cm}{\textbf{C)}} & \grafpos{1}{5}{2}{K}{1} &
\raisebox{1.9cm}{\textbf{D)}} & \grafpos{7}{2}{2}{K}{7}
\end{tabular}
\end{center}

%---------------------------------------
% ITEM 6
%---------------------------------------
\Needspace{0.84\textheight}
\item
A continuaci\'on, se muestra la representaci\'on gr\'afica, en la que las unidades est\'an
en kil\'ometros, de la ubicaci\'on ($P$) de un puerto, ($R$) de un barco y de la
circunferencia que corresponde al alcance m\'aximo de la se\~nal que emite el radar de
ese barco.

\begin{center}
\begin{tikzpicture}[scale=0.65]
  \path[use as bounding box] (-1.5,-1) rectangle (8,7);
  \draw[<->, thick] (-1,0) -- (7.5,0);
  \draw[<->, thick] (0,-0.5) -- (0,6.5);
  \node[anchor=west] at (7.7,-0.3) {$x$ (este)};
  \node[anchor=south east] at (1.1,6.7) {$y$ (norte)};
  \fill (0,0) circle (0.18) node[below left, inner sep=2pt] {$P$};
  \draw[thick] (5,3) circle (1);
  \draw[dashed] (5,0) -- (5,4);
  \draw[dashed] (0,3) -- (5,3);
  \draw[dashed] (0,4) -- (5,4);
  \fill (5,3) circle (0.18)
    node[anchor=south west, inner sep=2pt] {$R$};
  \fill (5,4) circle (0.18);
  \node[anchor=north] at (5,0) {$5$};
  \node[left] at (0,3) {$3$};
  \node[left] at (0,4) {$4$};
\end{tikzpicture}
\end{center}

De acuerdo con la informaci\'on anterior, si dos horas despu\'es la nueva ubicaci\'on ($S$)
del barco es $2$ km al norte de ($R$), entonces \textquestiondown cu\'al es la representaci\'on
gr\'afica, en la que las unidades est\'an en kil\'ometros, de la circunferencia que corresponde
al alcance m\'aximo de la se\~nal que emite el radar del barco en su nueva ubicaci\'on?

\vspace{1.5\baselineskip}
\begin{center}
\setlength{\tabcolsep}{0pt}
\renewcommand{\arraystretch}{1}
\begin{tabular}{@{}r@{\hspace{0.4cm}}c@{\hspace{1.4cm}}r@{\hspace{0.4cm}}c@{}}
\raisebox{1.9cm}{\textbf{A)}} & \grafpos{5}{5}{1}{S}{5} &
\raisebox{1.9cm}{\textbf{B)}} & \grafpos{7}{3}{1}{S}{7} \\[0.6cm]
\raisebox{1.9cm}{\textbf{C)}} & \grafpos{5}{1}{1}{S}{5} &
\raisebox{1.9cm}{\textbf{D)}} & \grafpos{3}{3}{1}{S}{3}
\end{tabular}
\end{center}

%---------------------------------------
% ITEM 7
%---------------------------------------
\Needspace{0.82\textheight}
\item
Un dispositivo electr\'onico, que emite una se\~nal inal\'ambrica con forma circular, se
ubica $8$ m al este y $4$ m al norte del laboratorio de una escuela. El alcance m\'aximo de
la se\~nal emitida por ese dispositivo es de $4$ m a su alrededor.

La siguiente representaci\'on gr\'afica, en la que las unidades est\'an en metros, muestra
las ubicaciones $V$ del dispositivo, $L$ del laboratorio y $C$ de la circunferencia
correspondiente al alcance m\'aximo de la se\~nal emitida por ese dispositivo\text{:}

\begin{center}
\begin{tikzpicture}[scale=0.46]
  \path[use as bounding box] (-1.5,-1) rectangle (15,11);
  \draw[<->, thick] (-1,0) -- (14,0);
  \draw[<->, thick] (0,-0.5) -- (0,10.5);
  \node[anchor=west] at (14.2,-0.3) {$x$ (este)};
  \node[anchor=south east] at (1.2,10.7) {$y$ (norte)};
  \fill (0,0) circle (0.175) node[below left, inner sep=2pt] {$L$};
  \draw[thick] (8,4) circle (4);
  \draw[dashed] (8,0) -- (8,4);
  \draw[dashed] (0,4) -- (8,4);
  \draw[dashed] (8,4) -- (8,8);
  \draw[dashed] (0,8) -- (8,8);
  \fill (8,4) circle (0.175)
    node[anchor=north east, inner sep=2pt, xshift=2pt, yshift=-2pt] {$V$};
  \fill (8,8) circle (0.175) node[above right, inner sep=2pt] {$C$};
  \node[anchor=north] at (8,0) {$8$};
  \node[left] at (0,4) {$4$};
  \node[left] at (0,8) {$8$};
\end{tikzpicture}
\end{center}

De acuerdo con la informaci\'on anterior, si desde su ubicaci\'on actual el dispositivo se
traslada $6$ m al oeste y $1$ m al sur, entonces \textquestiondown cu\'al es la
representaci\'on algebraica, en la que las unidades est\'an en metros, de la circunferencia
correspondiente al alcance m\'aximo de la se\~nal emitida por el dispositivo en su nueva
ubicaci\'on?

\begin{enumerate}[label=\Alph*)]
  \opcion{A}{(x-2)^{2}+(y-3)^{2}=4}
  \opcion{B}{(x-2)^{2}+(y-5)^{2}=16}
  \opcion{C}{(x-2)^{2}+(y-3)^{2}=16}
  \opcion{D}{(x-14)^{2}+(y-3)^{2}=16}
\end{enumerate}

\end{enumerate}

\vspace{1.2cm}
\begin{center}
\begin{tabular}{|c|c|c|c|c|c|c|c|}
\hline
\textbf{\'Item} & \textbf{1} & \textbf{2} & \textbf{3} & \textbf{4} &
                  \textbf{5} & \textbf{6} & \textbf{7} \\
\hline
\textbf{Clave}  & A & B & C & A & D & A & C \\
\hline
\end{tabular}
\end{center}

\end{document}
